% =======================================================
% 3. ARCHITETTURA DEL SISTEMA
% =======================================================
\section{Architettura del Sistema}

\subsection{Architettura Generale}
Il sistema \textit{Second Brain} è basato su un'architettura di tipo Client-Server per garantire la separazione delle responsabilità (Separation of Concerns) tra la gestione dell'interfaccia utente (editor Markdown) e l'orchestrazione delle chiamate ai Modelli di Linguaggio di Grande Scala (LLM), in linea con il vincolo R1-V-O.

Il Frontend Vue.js gestisce l'intera interazione con l'utente e non comunica mai direttamente con il servizio LLM: ogni richiesta di elaborazione AI transita dal Backend FastAPI, che funge da unico punto di orchestrazione verso il provider esterno. Questa scelta, oltre a rispettare il vincolo R2-V-O sullo standard delle chiamate API, evita di esporre nel browser eventuali credenziali di accesso al servizio LLM.

\begin{figure}[H]
	\centering
	\includegraphics[width=0.85\textwidth]{deployment.png}
	\caption{Architettura di alto livello e di deployment del sistema.}
	\label{fig:arch-generale}
\end{figure}

\subsection{Architettura di Deployment}
\label{sec:arch-deployment}
Frontend e Backend sono eseguiti in due container Docker distinti, orchestrati tramite Docker Compose (Figura~\ref{fig:arch-generale}):
\begin{itemize}
	\item \textbf{frontend}: applicazione Vue.js, servita in sviluppo dal dev server Vite e in produzione da un web server statico;
	\item \textbf{backend}: servizio FastAPI eseguito tramite Uvicorn, espone le API REST consumate dal frontend e contatta il servizio LLM esterno tramite HTTPS, con la relativa API key iniettata come variabile d'ambiente.
\end{itemize}
Non è presente alcun container di persistenza: come descritto in Sezione~\ref{sec:persistenza}, il gruppo ha deciso di non implementare la persistenza remota, funzionalità opzionale del capitolato (R4-V-Op, R78--R86); tutti i dati sono gestiti tramite il file system locale del browser, senza alcun componente server-side stateful.

Coerentemente con quanto indicato in \link{link_analisi_stato_arte}{Analisi dello Stato dell'Arte}, Sezione 4, la separazione in container non costituisce di per sé una scelta architetturale definitiva quanto piuttosto uno strumento per rendere l'ambiente di esecuzione riproducibile, portabile e semplice da avviare con un unico comando.

\subsection{Architettura del Frontend}
\label{sec:arch-frontend}
Il Frontend gestisce l'interazione dell'utente con l'editor Markdown e la visualizzazione dello stato delle operazioni AI. Il pattern architetturale adottato è \textbf{Model-View-Controller} nella forma basata su Observer: il Model realizza un'interfaccia \texttt{Subject} (Observable); la View realizza sia \texttt{Subject} sia \texttt{Observer}, potendo quindi essere osservata a sua volta dal Controller; il Controller realizza \texttt{Observer}. Il Controller osserva la View, la View osserva il Model; il Controller mantiene inoltre un riferimento diretto al Model per invocare comandi, senza però riceverne le notifiche direttamente.

Questa separazione esplicita, indipendente dal framework di rendering, mantiene la logica di dominio (Model e Controller) testabile in isolamento: i componenti Vue.js costituiscono l'implementazione concreta della View, limitandosi a riflettere nel template lo stato esposto dal relativo oggetto View e a inoltrare gli eventi dell'utente, senza contenere essi stessi logica applicativa.

\begin{figure}[H]
	\centering
	\includegraphics[width=0.8\textwidth]{pkg_frontend.png}
	\caption{Diagramma dei package del Frontend.}
	\label{fig:pkg-frontend}
\end{figure}

\subsubsection{Gestione dello Stato e dell'Editor}
Il \texttt{NoteModel} mantiene lo stato osservabile della nota corrente (contenuto, identificativo, flag di modifica non salvata) e delega la manipolazione testuale a un \texttt{MarkdownContentEditor}. Ogni modifica (formattazione, tabelle, elenchi, link, inserimento di testo) è incapsulata in un comando (Sezione~\ref{sec:design-pattern}) eseguito tramite \texttt{NoteModel.executeCommand()}, che aggiorna una \texttt{CommandHistory} a due pile per supportare Undo/Redo (R43-F-O, R44-F-O) in modo uniforme per tutte le operazioni di editing. Il salvataggio e l'apertura della nota (R75-F-O, R76-F-O) sono delegati tramite l'interfaccia \texttt{NoteService} a un'implementazione basata sul File System locale, così da non vincolare il Model a un meccanismo di persistenza specifico.

\subsubsection{Gestione Interazioni AI (Widget/Pannello AI)}
L'\texttt{AIRequestModel} mantiene lo stato della richiesta AI corrente come union type (\texttt{Idle}, \texttt{Processing}, \texttt{ProposalReady}, \texttt{Error}), riflesso dal componente \texttt{AIPanelView} tramite il meccanismo Observer (R72-F-O: visualizzazione continua dello stato di avanzamento). L'utente può accettare, rifiutare, rigenerare o interrompere manualmente una proposta AI (R69--R71, R73-F-O): l'accettazione riusa lo stesso meccanismo di comando dell'editing manuale, inserendo il testo proposto tramite un comando eseguito su \texttt{NoteModel} (Sezione~\ref{sec:design-pattern}, pattern Command). Le richieste verso il Backend sono mediate dall'interfaccia \texttt{AIService}, la cui implementazione concreta (Proxy) incapsula i dettagli della chiamata REST.

\begin{figure}[H]
	\centering
	\includegraphics[width=\textwidth]{class_frontend_model.png}
	\caption{Diagramma delle classi del livello Model del Frontend.}
	\label{fig:class-frontend-model}
\end{figure}

\begin{figure}[H]
	\centering
	\includegraphics[width=\textwidth]{class_frontend_vcp.png}
	\caption{Diagramma delle classi dei livelli View, Controller e Proxy del Frontend.}
	\label{fig:class-frontend-vcp}
\end{figure}

\subsection{Architettura del Backend}
Il Backend funge da orchestratore centrale e da strato di astrazione verso i provider LLM esterni. È organizzato in due package Python distinti: il livello applicativo/di presentazione (router REST, DTO, Dependency Injection) e il dominio AI (\texttt{AI\_Domain}), che contiene la logica di costruzione dei prompt e di interazione con l'LLM indipendente dal framework web.

\begin{figure}[H]
	\centering
	\includegraphics[width=0.9\textwidth]{pkg_backend.png}
	\caption{Diagramma dei package del Backend.}
	\label{fig:pkg-backend}
\end{figure}

\subsubsection{Livello di Presentazione / Controller}
I router FastAPI (\texttt{AIRouter}, \texttt{ExportRouter}) ricevono le richieste dal client, le validano tramite i DTO di \texttt{api/schemas} (\texttt{AIOperationRequest}, \texttt{ExportRequest}) e delegano l'elaborazione ai servizi applicativi, ottenuti tramite Dependency Injection (\texttt{DIProviders}). Un gestore di errori centralizzato (\texttt{ErrorHandlers}) traduce le eccezioni di dominio in risposte HTTP coerenti (Sezione~\ref{sec:gestione-errori}).

\begin{figure}[H]
	\centering
	\includegraphics[width=0.85\textwidth]{class_backend_api_layer.png}
	\caption{Diagramma delle classi del livello di presentazione e DI del Backend.}
	\label{fig:class-backend-api}
\end{figure}

\subsubsection{Livello di Business Logic (Core LLM)}
L'\texttt{AIService} è il motore centrale del dominio AI: riceve il tipo di operazione richiesta (riassunto, traduzione, riscrittura, Distant Writing, o una delle sei prospettive dei ``Cappelli di De Bono''), la risolve tramite un'\texttt{AIOperationFactory} auto-registrante, e ne invoca \texttt{build\_prompt()} in modo polimorfico. Ogni operazione costruisce il proprio \texttt{Prompt} tramite un \texttt{PromptBuilder} fluente, isolando la logica specifica di ciascuna trasformazione testuale (Sezione~\ref{sec:design-pattern}, pattern Strategy, Factory Method e Builder).

\begin{figure}[H]
	\centering
	\includegraphics[width=\textwidth]{class_ai_domain_operations.png}
	\caption{Diagramma delle classi: operazioni AI e costruzione dei prompt (\texttt{AI\_Domain/domain}).}
	\label{fig:class-ai-domain-ops}
\end{figure}

\subsubsection{Livello di Integrazione / Infrastructure}
Il pacchetto \texttt{AI\_Domain/llm} maschera le specificità del provider LLM concreto dietro l'interfaccia \texttt{LLMService}. Una catena di decorator (\texttt{CachingLLMAdapter}, \texttt{LoggingLLMAdapter}) aggiunge funzionalità trasversali di caching e logging senza modificare l'implementazione di base; l'ultimo anello della catena è un adapter verso l'API compatibile con lo standard OpenAI (modelli Gemma, ChatGPT-OSS o Qwen, R2-V-O, R3-V-O), configurabile tramite \texttt{base\_url} per puntare all'endpoint messo a disposizione dall'azienda proponente.

\begin{figure}[H]
	\centering
	\includegraphics[width=0.9\textwidth]{class_ai_domain_llm.png}
	\caption{Diagramma delle classi: integrazione con il servizio LLM (\texttt{AI\_Domain/llm}).}
	\label{fig:class-ai-domain-llm}
\end{figure}

\FloatBarrier
